%%%%%%%%%%%%%%%%%%%%%%%%%%%%%%%%
%% Printout-Cards
%%%%%%%%%%%%%%%%%%%%%%%%%%%%%%%%

\ExplSyntaxOn

\NewDocumentEnvironment{RpgCard}{ O{} +b }
{
	%% Do this before group because it does set a global flag

	\group_begin:
	\__rpg_card_EntryCheck:
	\RpgCardSet{#1}

	\tl_gclear:N \__rpg_card_fontState
	\tl_gclear:N \__rpg_card_footnote
	\RenewDocumentCommand{\cardbreak}{}{ \par \penalty -10000 \mbox{} } %% HACK: mbox or vspace?
	\RenewDocumentCommand{\footnote}{m}{ \__rpg_card_footnoteChar\tl_gset:Nn{\__rpg_card_footnote}{##1} }


	\__rpg_card_processContents:n{#2}
}
{
	\group_end:
}

%%%%%%%%%%%%%%%
% Options
%%%%%%%%%%%%%%%


%% For a single-source-of-truth on defaults, we leave them blank here, and set them in the Reset command
\RpgThemeKeys{ rpg / card}
{
	\RpgThemeTLKey{width}{\__rpg_card_width}{}
	\RpgThemeTLKey{height}{\__rpg_card_height}{}
	\RpgThemeTLKey{hmargin}{\__rpg_card_hmargin}{}
	\RpgThemeTLKey{vmargin}{\__rpg_card_vmargin}{}
	\RpgThemeTLKey{cardsep}{\__rpg_card_cardsep}{}
	\RpgThemeTLKey{footnote}{\__rpg_card_footnoteChar}{*}

	\RpgThemeBoolKey{vcentre}{\__rpg_card_vcentre}
	\RpgThemeBoolKey{vcenter}{\__rpg_card_vcentre}
	%% appearance
	\RpgThemeTLKey{color}{\__rpg_card_color}{}
	\RpgThemeTLKey{opacity}{\__rpg_card_opacity}{}
}

% \begin{noindent}
\NewDocumentCommand\RpgCardReset{}
{
	\tcbset{rpgcard-composite/.style={rpgcard}} %% resets the composite back to the baseline
	\RpgCardSet{
		width=6.75cm,
		height=9.7cm,
		hmargin=0.2cm,
		vmargin=0.2cm,
		color=white,
		cardsep=5pt,
		opacity=1,
	}
}
\tcbset{rpgcard/.style={
			valign=center,
			halign=left,
			colframe=black,
			size=tight,
			arc=1mm,
			auto~outer~arc,
			enhanced~jigsaw,
		}}

\NewDocumentCommand\RpgCardSet{+m}
{
	%%intercept the known values
	\keys_set_known:nnN {rpg/card}{#1}{\l__tmpa_tl}

	%%insert the rest in a specific order to ensure the correct resolution (color overrides colback etc)
	%%Rule out titles from being inserted (they break the card splitting rules)
	\exp_args:Nx \tcbset
	{
		rpgcard/.append~style={
		\l__tmpa_tl, 
		title={},
		colback=\__rpg_card_color,
		opacityback=\__rpg_card_opacity,
		width=\__rpg_card_width,
		height=\__rpg_card_height,
		} 
	}


	%%% now set the text box size 

	\dim_set:Nn \__rpg_card_textWidth{
		\dim_eval:n { \__rpg_card_width - (\__rpg_card_hmargin * 2) }
	}
	\dim_set:Nn \__rpg_card_textHeight{
		\dim_eval:n { \__rpg_card_height - (\__rpg_card_vmargin * 2) }
	}

	%% Now check nothing went negative
	\dim_compare:nT{\__rpg_card_textWidth <= 0pt}
	{
		\msg_error:nneeee {rpg}{card-dimensions}{hmargin}{\__rpg_card_hmargin}{\dim_eval:n{\__rpg_card_width /2}}{width}
	}
	\dim_compare:nT{\__rpg_card_textHeight <= 0pt}
	{
		\msg_error:nneeee {rpg}{card-dimensions}{vmargin}{\__rpg_card_vmargin}{\dim_eval:n{\__rpg_card_height /2}}{height}
	}

}

%\end{noindent}

\tcbset{rpgcard-composite/.style={}}
\newtcolorbox{__RpgCard}[1][]{
	rpgcard,
	#1
}


%%%%%%%%%%%%%%%%%
%% Helpers
%%%%%%%%%%%%%%%%%
\bool_new:N \__rpg_card_isNested
\bool_set_false:N \__rpg_card_isNested
\msg_new:nnn {rpg}{card-nested}{Cannot~nest~RpgCard~inside~RpgCard}
\cs_new_protected:Npn \__rpg_card_EntryCheck:{
	\bool_if:NT \__rpg_card_isNested
	{
		\msg_error:nnn{rpg}{card-nested}{}
	}
	\bool_set_true:N \__rpg_card_isNested
}


%%%%%%%%%%%%%%%%%
% Page splitting
%%%%%%%%%%%%%%%%%

\NewDocumentCommand{\cardbreak}{}{} %% command does nothing until redefined in an RpgCard env.
\dim_new:N \__rpg_card_textWidth
\dim_new:N \__rpg_card_textHeight
\msg_new:nnn { rpg } { card-dimensions}{~\\Dimension~error!\\The~#1~is~too~large~(\dim_eval:n{#2}).~It~cannot~exceed~#3~(half~the~#4).}
\bool_new:N \__rpg_card_isFirstPage
\box_new:N \__rpg_card_contents
\box_new:N \__rpg_card_subContents
\box_new:N \__rpg_card_tmpbox

\cs_new_protected:Npn \__rpg_card_processContents:n#1
{
	%% places the contents of the card into a mega-card (\__rpg_card_contents)
	\__rpg_card_buildBoundingBox:n{#1}

	\ignorespaces\noindent\leavevmode

	%% now we do the iteration
	\bool_set_true:N \__rpg_card_isFirstPage
	\bool_do_until:nn {\box_if_empty_p:N \__rpg_card_contents}
	{
		\splittopskip=0pt

		%%split off a chunk of the megabox to be our single-card
		\vbox_set_split_to_ht:NNn \__rpg_card_subContents \__rpg_card_contents { \__rpg_card_textHeight }
		\bool_if:NTF \__rpg_card_isFirstPage
		{
			\bool_set_false:N \__rpg_card_isFirstPage
		}
		{
			\hskip\__rpg_card_cardsep
		}

		\bool_if:NTF \__rpg_card_vcentre
		{
			\def\insert{}
			\vbox_set:Nn \__rpg_card_tmpbox
			{
				\unvbox \__rpg_card_subContents
				\unskip
				\unkern
			}

			% Repack into fixed target height with centering glues
			\vbox_set_to_ht:Nnn \__rpg_card_subContents { \__rpg_card_textHeight }
			{
				\vfill
				\box_use_drop:N \__rpg_card_tmpbox
				\vfill
			}
		}
		{
			\def\insert{\vspace{\__rpg_card_vmargin}}
		}
		\hbox{
			\begin{__RpgCard}
				%% insert the previous font state
				\__rpg_card_fontState

				\hfil {{\box_use:N \__rpg_card_subContents} }\hfil

				%% capture the current font state
				\tl_gset:Nn \__rpg_card_fontState{\the\font}

			\end{__RpgCard}
		}

	}



}
\cs_new_protected:Npn \__rpg_card_buildBoundingBox:n#1
{
	\vbox_set:Nn \__rpg_card_contents{
		\hsize = \__rpg_card_textWidth
		\linewidth = \hsize

		\RpgFontCardBody

		#1


		\tl_if_empty:eF{\__rpg_card_footnote}
		{
			\vfil
			\vbox{
				\hrule
				\vspace{0.5em}

				\noindent \__rpg_card_footnoteChar~\__rpg_card_footnote
			}
		}

	}
}



%%%%% And at the very end....

\RpgCardReset{} %(place it into a default state after everything has been declared)

